\cotefile{29}
\scanpage{25}
% archive sheet number: 24
évidente lorsque $P=G'_S$ avec l'opération régulière droite de $G'$, et l'opération gauche de $G$ est donnée par un monomorphisme $G\quad G'$; on se ramène à ce cas particulier par localisation étale.

On peut se proposer de normaliser la notion de donnée de ramification, pour éliminer l'arbitraire dans le choix d'un $\underline{R}$. Mais alors \unsure{\deleted{o}il|\deleted{i}ol|\deleted{b}il|\deleted{d}il|\deleted{i}bl|\deleted{i}dl} semble qu'on soit infailliblement conduit à prendre comme catégorie $\underline{C}$ la catégorie $\underline{\mathrm{Rev}}^{\underline{R}}$ elle-même qu'on se propose \unsure{d\deleted{'}e|d\deleted{e}'} construire et d'étudier ! (Comparer avec la situation analogue pour les relations entre sites et topos, l'objet auxiliaire étant le site, l'objet intrinsèque de nature géométriquenétant le topos associé à ce site).

2. \uline{Propriétés de la catégorie} $\operatorname{Rev}^{\underline{R}}(S)$. \uline{Le groupe fondamental} $ {}^{\underline{R}}_1(S,\quad)$.

On constate immédiatement que la catégorie $\operatorname{Rev}^{\underline{R}}(S)$ hérite de toutes les propriétés d'ex\unsure{a|\deleted{e}a|\deleted{a}e}\unsure{s|c|\deleted{s}c|\deleted{c}s}titude \deleted{déjà} établies pour la catégorie des revêtementsétales d'un schémas donné, et dont une partie sont énoncés dans Exp V n\textsuperscript{\uline{o}}4 (conditions G1 à G3). De façon précise, on peut dire que dans la catégorie $\operatorname{Rev}^{\underline{R}}(S)$, les limites projectives finies et les limites inductives finies existent, et satisfont aux mêmes propriétés d'exactitude que dans la catégorie des ensembles (telle l'exactitude du foncteur changement de base dans cette catégorie); \deleted{de plus\unsure{;|,}} \inserted{en particulier} tout morphisme de la catégorie se factorise en le produit d'un épimorphisme et d'un mnomorphisme, \inserted{les épimorphismes et monomorphismes sont \deleted{\unsure{s}tri} effectifs; enfin} tout monomorphisme est un isomorphisme avec un sommande direct. On fera attention que le foncteur ``réalisation géométrique'' sur $\operatorname{Rev}^{\underline{R}}(S)$ commute bien aux limites inductives finies (et \deleted{par suite au\unsure{x}} en particulier aux sommes finies), mais non pas en général aux limites projectives finies, et en particulier, pas aux produits de deux objets.

Pour appliquer la théorie de Galois générale de V 4, il faut dis-

